\section*{الفصل \ref{ch:g10:drainage-basins} --- أحواض التصريف والأنهار والفيضانات}

\begin{solution}{exo:g10:drainage-basins:1}
\omterm{def:g10:drainage-basins:drainage-basin}{حوض التصريف} لنقطة من نهر هو المساحة كلها من الأرض التي يجري ماؤها السطحي إلى
النهر في أعلى تلك النقطة ويمر بها. \omterm{def:g10:drainage-basins:drainage-basin}{وخط تقسيم المياه} هو الخط، الممتد على القمم، الذي
يفصل بين حوضين متجاورين. والقطرتان الساقطتان على بعد متر واحدة من الأخرى على جانبي
خط تقسيم تنحدران على سفحين متقابلين، إلى حوضين مختلفين: وقد يبلغ نهراهما بحرين
مختلفين.
\end{solution}

\begin{solution}{exo:g10:drainage-basins:2}
\omterm{def:g10:drainage-basins:discharge}{التصريف} هو حجم الماء الذي يعبر مقطعًا عرضيًا للنهر في كل ثانية، بالوحدة
\unit{m^3/s}. و\qty{1}{m^3/s} يساوي \qty{1000}{L} في كل ثانية؛ وفي اليوم،
$1000 \times 86\,400 = 86.4$ مليون لتر.
\end{solution}

\begin{solution}{exo:g10:drainage-basins:3}
قبل الفيضان حمل النهر نحو \qty{330}{m^3/s} (بين نحو
\num{240} و\qty{390}{m^3/s}). وكانت الذروة \qty{9231}{m^3/s} في
21 يناير. وبدأ الارتفاع في 18 يناير (\qty{362}{m^3/s}): فقد استغرق ثلاثة
أيام.
\end{solution}

\begin{solution}{exo:g10:drainage-basins:4}
$1/50 = 0.02$، أي \qty{2}{\percent}. لا: فالسنوات مستقلة؛ والفرصة
\qty{2}{\percent} من جديد هذه السنة، مهما حدث في السنة الماضية.
\end{solution}

\begin{solution}{exo:g10:drainage-basins:5}
القدم المكعب $0.3048^3 \approx \qty{0.02832}{m^3}$، إذن
$593\,000 \times 0.02832 \approx \qty{16800}{m^3/s}$. فالأمازون يحمل
$206\,000/16\,800 \approx 12$ ضعفًا من ذلك.
\end{solution}

\begin{solution}{exo:g10:drainage-basins:6}
مساحات الشرائح: $1.5 \times 0.3 = 0.45$ و$0.90$ و$1.05$ و\qty{0.60}{m^2}:
$A = \qty{3.0}{m^2}$. والتصريفات: $0.2 \times 0.45 = 0.09$ و$0.5 \times 0.90
= 0.45$ و$0.6 \times 1.05 = 0.63$ و$0.3 \times 0.60 = 0.18$: $Q =
\qty{1.35}{m^3/s}$. والسرعة المتوسطة $1.35/3.0 = \qty{0.45}{m/s}$.
\end{solution}

\begin{solution}{exo:g10:drainage-basins:7}
$66 \times 10^6 \times 0.020 = \qty{1.32e6}{m^3}$. واليومان
$2 \times 86\,400 = \qty{172800}{s}$: $1.32 \times 10^6/172\,800 \approx
\qty{7.6}{m^3/s}$ في المتوسط.
\end{solution}

\begin{solution}{exo:g10:drainage-basins:8}
مرتبةً: 120، 88، 65، 59، 51، 47، 42، 38، 35. ومع $n = 9$، $T = 10/m$: فللأكبر
(120) $T = 10$ سنوات، وللثانية (88) 5 سنوات، وللثالثة (65)
نحو 3.3 سنة. والوسيطة، 51، رتبتها 5 و$T = 2$ سنة: فهي تُساوى
أو تُتجاوز سنة من كل سنتين.
\end{solution}

\begin{solution}{exo:g10:drainage-basins:9}
\omterm{def:g10:drainage-basins:hydrograph}{هيدروغراف} البلدة يرتفع أبكر وأعلى بكثير، وينخفض بسرعة: فعلى السفوح الشديدة الانحدار
والأسطح والطرق لا يتسرب إلا قليل من الماء ويبلغ الجريان السطحي الجدول بسرعة ودفعة
واحدة؛ وزمن تأخره قصير. أما \omterm{def:g10:drainage-basins:hydrograph}{هيدروغراف} الحوض المكسو بالغابات فيرتفع متأخرًا، إلى ذروة
أدنى وأعرض: إذ تمتص التربة والغابة قدرًا كبيرًا من المطر وتطلقانه ببطء، ويجري الماء
ببطء على السفوح اللطيفة. ويبقى تدفقه القاعدي أعلى بعد العاصفة.
\end{solution}

\begin{solution}{exo:g10:drainage-basins:10}
$9231 \times 86\,400 \approx \qty{7.98e8}{m^3}$، موزعة على
$29\,900 \times 10^6 = \qty{2.99e10}{m^2}$: أي نحو \qty{0.027}{m}، أو
\qty{27}{mm}. \omterm{def:g10:drainage-basins:hydrograph}{والتدفق القاعدي} يعطي $330 \times 86\,400/2.99 \times 10^{10}
\approx \qty{0.001}{m}$، أي نحو \qty{1}{mm} في اليوم: فقد حمل يوم الذروة نحو
28 ضعفًا.
\end{solution}

\begin{solution}{exo:g10:drainage-basins:11}
$P = 1 - 0.9^{10} \approx 1 - 0.349 = 0.651$: أي نحو \qty{65}{\percent}. ولا
يأتي إطلاقًا باحتمال $0.9^{10} \approx 0.35$، أي نحو
\qty{35}{\percent}.
\end{solution}

\begin{solution}{exo:g10:drainage-basins:12}
\omterm{def:g10:drainage-basins:hydrograph}{الهيدروغراف} مثلث طوله ثلاثة أيام وارتفاعه \qty{1200}{m^3/s}: مساحته
$\tfrac12 \times 1200 \times 3 \times 86\,400 \approx
\qty{1.56e8}{m^3}$، أي \num{156} مليون متر مكعب. والجزء الواقع فوق
\qty{800}{m^3/s} مثلث أصغر له الشكل نفسه، أصغر بثلاث مرات في الارتفاع
($400$ بدلًا من $1200$) وفي المدة (يوم واحد بدلًا من
ثلاثة): $\tfrac12 \times 400 \times 86\,400 \approx \qty{1.73e7}{m^3}$،
أي \num{17} مليون متر مكعب. وهذا أقل من \num{60} مليونًا يستطيع
الخزان تخزينها: نعم.
\end{solution}

\begin{solution}{exo:g10:drainage-basins:13}
حين يمتلئ النهر حتى ضفتيه يكون عمقه $800/(100 \times 2) = \qty{4}{m}$.
(1) يحمل السهل الفيضي \qty{400}{m^3/s} على عرض \qty{900}{m} بسرعة
\qty{0.5}{m/s}: فعمقه $400/(900 \times 0.5) \approx \qty{0.9}{m}$.
(2) بين الحواجز الترابية: عمق $1200/(100 \times 2) = \qty{6}{m}$، أي \qty{2}{m}
فوق ضفتيه، وهذا أكثر من ضعف العمق على السهل الفيضي دون
حواجز.
\end{solution}

\begin{solution}{exo:g10:drainage-basins:14}
نريد $(1 - 1/T)^{30} = 0.5$، إذن $1 - 1/T = 0.5^{1/30} \approx 0.97716$،
و$1/T \approx 0.02284$ و$T \approx 44$ سنة.
\end{solution}

\begin{solution}{exo:g10:drainage-basins:15}
يرسل كل مربع ماءه إلى أدنى جيرانه (الصفوف من الشمال إلى الجنوب،
والغرب في اليسار والشرق في اليمين):
\begin{center}
\begin{tabular}{|c|c|c|c|}
\hline
غرب & غرب & غرب & مخرج \\ \hline
جنوب & غرب & غرب & شمال \\ \hline
جنوب & جنوب & شمال & شمال \\ \hline
مخرج & شرق & شرق & شمال \\ \hline
\end{tabular}
\end{center}
باتباع الأسهم، تصرف المربعات 47 و39 و44 و37 و46 و30 ماءها إلى المخرج
الجنوبي الشرقي: 6 مربعات، أي \qty{0.06}{km^2}. والمربعات العشرة الأخرى تصرف ماءها
إلى المخرج الشمالي الغربي: \qty{0.10}{km^2}. ويمتد خط التقسيم من الحافة
الشرقية بين 50 و47، ثم بين 48 و47، وبين 48 و44، وبين
49 و44، وبين 49 و46، وبين 51 و46 حتى الحافة الجنوبية.
\end{solution}

\begin{solution}{pb:g10:drainage-basins:1}
\textbf{1.} $11\,560 \times 1.609^2 \approx \qty{29900}{km^2}$.

\textbf{2.} \omterm{def:g10:drainage-basins:drainage-basin}{خط تقسيم المياه} للحوض الواقع في أعلى المحطة.

\textbf{3.} $484\,000 \times 0.3048^3 \approx \qty{13700}{m^3/s}$.

\textbf{4.} $13\,705/3341 \approx 4.1$ مرة.

\textbf{5.} $A = 400 \times 8 = \qty{3200}{m^2}$ و$v = Q/A = 13\,705/3200
\approx \qty{4.3}{m/s}$.

\textbf{6.} نحو \qty{330}{m^3/s}.

\textbf{7.} \qty{9231}{m^3/s}، في 21 يناير.

\textbf{8.} ثلاثة أيام.

\textbf{9.} $9231 \times 86\,400/(29\,900 \times 10^6) \approx
\qty{0.027}{m}$: أي نحو \qty{27}{mm}.

\textbf{10.} بلغ النهرَ ماء جديد: فقد أرسل هطول آخر على الحوض (أو ثلج ذائب)
نبضة ثانية أصغر من الجريان السطحي.

\textbf{11.} بعد الذروة بسبعة عشر يومًا. فالارتفاع يأتي من الجريان السطحي للسفوح
القريبة من الجداول، الذي يصل بسرعة ودفعة واحدة؛ أما على الطرف الهابط فما زال النهر
يتلقى ماء أبعد أجزاء الحوض والماء الذي ينصرف ببطء من التربة.

\textbf{12.} $T = n + 1 = 96$ سنة.

\textbf{13.} $96/2 = 48$ سنة.

\textbf{14.} $96/3 = 32$ سنة.

\textbf{15.} $96/4 = 24$ سنة و$96/5 = 19.2$ سنة. لا يهم كثيرًا: ففي الحالتين
يعود هذا \omterm{def:g10:drainage-basins:discharge}{التصريف} نحو مرة كل عشرين سنة.

\textbf{16.} خمس سنوات (1936 و1937 و1943 و1972 و1996): $5/95 \approx
0.053$، أي نحو \qty{5}{\percent} في السنة، \omterm{def:g10:drainage-basins:return-period}{وفترة عودة} قدرها نحو
$96/5 \approx 19$ سنة.

\textbf{17.} $1/100 = 0.01$، أي \qty{1}{\percent}.

\textbf{18.} $1 - 0.01 = 0.99$.

\textbf{19.} $0.99^{30} \approx 0.740$.

\textbf{20.} $1 - 0.740 = 0.26$: \textbf{فرصة حدوث فيضان مئة سنة واحد على الأقل
في 30 سنة نحو \qty{26}{\percent}}، أي واحدة من أربع.
\end{solution}
